\begin{theorem}[固定二维 \texorpdfstring{$2$}{2}-进 ambient 的充分性与有限秩同态账本的绝对不足]\label{thm:conclusion-prime-register-fixed-2adic-ambient-vs-finite-ledger}
\leanverified{paper\_conclusion\_prime\_register\_fixed\_2adic\_ambient\_vs\_finite\_ledger}
设 \(H_{\mathrm{can}}\) 为 canonical 有限时深 prime-register 子单子。则存在忠实单子表示
\[
\Psi:H_{\mathrm{can}}\hookrightarrow \operatorname{Aff}\!\bigl(T_2(E_{\min})\bigr),
\qquad
T_2(E_{\min})\cong \ZZ_2^2,
\]
亦即固定二维 \(2\)-进 Tate 模块已足以承载全部有限深度 canonical prime-register 动力学。

另一方面，不存在任何有限生成交换可消去加法账本 \(L\) 使得 cofinal 的乘法结构 \(\NN^\times\) 能忠实嵌入
\[
(\ZZ\oplus L,+)
\qquad\text{或}\qquad
(\NN^k,+)
\]
对任何有限 \(k\) 成立。因而 prime-multiplicative 机制的真正障碍不在 ambient 空间维数，而在要求把乘法语义保持为有限秩加法账本同态。
\end{theorem}
\begin{proof}
主文 H.160 已构造出 canonical 仿射作用
\[
\Psi(r,k)(x)=B^k x+\operatorname{ev}_B(r)
\]
并证明其在 \(H_{\mathrm{can}}\) 上忠实，其中底层空间正是
\[
T_2(E_{\min})\cong \ZZ_2^2.
\]
故固定二维 \(2\)-进 Tate 模块确已足够。

另一方面，主文 H.130--H.131 已证明
\[
\NN^\times\not\hookrightarrow (\NN^k,+)
\qquad(\forall\,k<\infty),
\]
且更强地，对任意有限生成交换群 \(L\) 都不存在忠实幺半群嵌入
\[
\NN^\times\hookrightarrow (\ZZ\oplus L,+).
\]
因此，有限维 ambient realization 的存在，与有限秩 additive ledger realization 的存在是两件根本不同的事；前者成立，后者失败。证毕。
\end{proof}

\begin{theorem}[一寄存器可序列化与线性同态债务的同时成立]\label{thm:conclusion-prime-register-one-register-serialization-vs-linear-hom-debt}
\leanverified{paper\_conclusion\_prime\_register\_one\_register\_serialization\_vs\_linear\_hom\_debt}
记
\[
S_r:=\{p_1,\dots,p_r\},
\qquad
M_r:=\NN^\times_{S_r}.
\]
则有精确维数差公式
\[
\operatorname{cdim}^{+}_{\mathrm{hom}}(M_r)=\frac r2,
\qquad
\operatorname{cdim}_{\mathrm{impl}}(M_r)=\frac12,
\qquad
\operatorname{cdim}^{+}_{\mathrm{hom}}(M_r)-\operatorname{cdim}_{\mathrm{impl}}(M_r)=\frac{r-1}{2}.
\]
换言之，同一对象族可以在实现层被无损序列化到单寄存器，而在结构层却需要线性增长的加法寄存器数才能保持幺半群语义。
\end{theorem}
\begin{proof}
主文已证明：任意单射幺半群同态
\[
\Phi_r:M_r\hookrightarrow (\NN^{k_r},+)
\]
必满足 \(k_r\ge r\)，并且该下界可达；因此
\[
\operatorname{cdim}^{+}_{\mathrm{hom}}(M_r)=\frac r2.
\]

另一方面，只要 \(r\ge 1\)，\(M_r\) 就是可数无限集。主文关于实现层半圆维的定义与其可数无限对象分类给出
\[
\operatorname{cdim}_{\mathrm{impl}}(M_r)=\frac12.
\]
更具体地，命题 13.4059 说明任意有限描述对象族都可通过可计算单射
\[
\mathrm{enc}:A^\ast\hookrightarrow \NN
\]
无损序列化到单寄存器；但该编码并不保持幺半群运算，因此不会降低结构层的同态圆维。

两式相减即得
\[
\operatorname{cdim}^{+}_{\mathrm{hom}}(M_r)-\operatorname{cdim}_{\mathrm{impl}}(M_r)=\frac{r-1}{2}.
\]
证毕。
\end{proof}
